\documentclass[10pt]{beamer}

% Packages
\usepackage{graphicx}
\usepackage{tikz}
\usepackage{bm}
\usepackage{xcolor}
\usepackage{pifont}
\usepackage{booktabs}
\usepackage{listofitems}

% tikz setup
\usetikzlibrary{calc}
\usetikzlibrary{decorations.pathreplacing}

% Metropolis options
\usetheme{metropolis}
\metroset{progressbar=frametitle,
          block=fill}

% notes environment: purple text
\newenvironment{notes}
    {\begingroup \color{violet}}
    {\endgroup}

\newcommand{\todo}[1]{{\color{red}[\textbf{TODO:} {#1}]}}

% maths environments
\theoremstyle{plain}\newtheorem{example}{Example}[section]
\theoremstyle{plain}\newtheorem{theorem}{Theorem}[section]
\theoremstyle{plain}\newtheorem{proposition}{Proposition}[section]
\theoremstyle{plain}\newtheorem{lemma}{Lemma}[section]
\theoremstyle{plain}\newtheorem{definition}{Definition}[section]
\theoremstyle{plain}\newtheorem{claim}{Claim}[section]
\theoremstyle{plain}\newtheorem{corollary}{Corollary}[section]

% preserve some default commands which we will overwrite
\let\sectionsymbol\S
\let\originalqedsymbol\qedsymbol

% proof environment for claims, which uses a different QED symbol
\newenvironment{claimproof}
    {\begin{proof}[Proof of claim.]\renewcommand{\qedsymbol}{$\diamondsuit$}}
    {\end{proof}\renewcommand{\qedsymbol}{\originalqedsymbol}}

% Environment for axiom definition
\newcommand{\axiomenvsep}{3mm}
\newcommand{\axiomenvmargin}{2.5em}
\newenvironment{axiom}[1][Axiom]%
    {\vspace{\axiomenvsep}\begin{adjustwidth}{\axiomenvmargin}{\axiomenvmargin}{#1}.}
    {\end{adjustwidth}\vspace{\axiomenvsep}}

% when defining several axioms in a row, we only want extra vertical space
% before the first and after the last. define a new environment which
% temporarily sets the spacing to 0
\newenvironment{axiomlist}%
    {\vspace{\axiomenvsep}\begingroup\renewcommand{\axiomenvsep}{0em}}
    {\endgroup\vspace{\axiomenvsep}}

% style for referencing axioms in the text
\newcommand{\axiomref}[1]{\textbf{#1}}

% inline list
\newlist{inlinelist}{enumerate*}{1}
\setlist[inlinelist]{label=(\roman*)}

% cleveref formats for equations and items
\crefformat{equation}{(#2#1#3)}
\Crefformat{equation}{(#2#1#3)}
\crefformat{enumi}{(#2#1#3)}
\Crefformat{enumi}{(#2#1#3)}
\crefformat{inlinelisti}{#2#1#3}
\Crefformat{inlinelisti}{#2#1#3}

% cleveref cannot deduce the plural form for the `claim` environment, so give
% it explicity
\crefname{claim}{Claim}{Claims}
\Crefname{claim}{Claim}{Claims}

% override the name for page number references
\Crefname{page}{p.}{p.}

% greek letters
\renewcommand{\phi}{\varphi}
\renewcommand{\epsilon}{\varepsilon}

% misc operators
\newcommand{\trcls}{+}  % transitive closure
\renewcommand{\cite}{\parencite}
\renewcommand{\cites}{\parencites}

% logical symbols
\newcommand{\limplies}{\rightarrow}
\newcommand{\liff}{\leftrightarrow}
\newcommand{\limpliedby}{\leftarrow}

\newcommand{\complexityclass}[1]{\textsf{#1}}


% Metadata
\title{
    Trustworthiness and Expertise: Social Choice and Logic-based Perspectives
}
\date{}
\author{
    Joseph Singleton
    \\singletonj1@cardiff.ac.uk
}
\titlegraphic{\includegraphics[width=1cm]{images/cardifflogo}}

\begin{document}

\maketitle

\begin{frame}{Thesis Overview}
    \begin{itemize}
        \item The thesis studies problems relating to \alert{unreliable
              information} and \alert{expertise}
        \item Emphasis on applying formal methods
        \begin{itemize}
            \item social choice theory
            \item modal logic
            \item belief revision
            \item formal learning theory
        \end{itemize}
    \end{itemize}
\end{frame}

\begin{frame}{Social Choice Perspectives}
    \begin{itemize}
        \item The first half of the thesis uses the methodology of
              computational social choice theory
          \item We develop an axiomatic framework for \alert{truth discovery}
              and \alert{bipartite tournament ranking}
        \item Axioms formalise ``desirable'' properties for a method
        \item Axiomatic analysis complements empirical work for comparing and
              developing new methods
    \end{itemize}
\end{frame}

\begin{frame}{Truth Discovery}
    \begin{itemize}
        \item Truth Discovery has recently arisen as a branch of the literature
              on crowdsourcing
        \item Central question: given conflicting information, who should we
              trust and what should we believe?
        \item We set out new axioms for truth discovery, and analyse an
              existing method from the literature
    \end{itemize}

    % Copied from thesis with minor modifications
    \begin{figure}
        \centering
        \begin{tikzpicture}[thick,scale=0.8]
            \LARGE
            \networkinitwithnames{s,t,u,v}{s,t,u,v}{c,d,e,f}{+,-,+,-}
            \object{o_1}{c}{d};
            \object{o_2}{e}{f};
            \report{s}{c};
            \report{s}{e};
            \report{t}{d};
            \report{t}{f};
            \report{u}{f};
            \report{v}{f};
        \end{tikzpicture}
    \end{figure}

\end{frame}

\begin{frame}{Bipartite Tournament Ranking}
    \begin{itemize}
        \item ``Ground truth'' data can help with truth discovery:
        \begin{itemize}
            \item We already know something about the trustworthiness of
                  sources
            \item But this is not straightforward if objects vary in
                  \emph{difficulty}
        \end{itemize}
        \item We generalise aspects of this problem: how should players in a
              \alert{bipartite tournament} be ranked?
    \end{itemize}

    \todo{Example tournament}
\end{frame}

\begin{frame}{Logic-based Perspectives}
    \begin{itemize}
        \item The second half of the thesis uses logic-based methods:
        \begin{itemize}
            \item Modal logic framework to reason about expertise
            \item Multi-source belief change problem with non-experts
            \item Investigation into truth-tracking with non-experts
        \end{itemize}
    \end{itemize}
\end{frame}

\begin{frame}{Logic of Expertise}
    \begin{nitemize}
        \item new logical framework with E and S modalities
        \item connections with epistemic logic
        \item investigation into different properties of expertise
    \end{nitemize}
\end{frame}

\begin{frame}{Multi-Source Belief Change}
    \begin{nitemize}
        \item belief change problem in the style of AGM and KP belief merging
        \item uses the expertise framework to reason about expertise and trust
            (we trust a source if we believe they are an expert)
        \item key feature: trust can be revised as well as beliefs about the
              world
        \item approach is axiomatic: rationality postulates
    \end{nitemize}
\end{frame}

\begin{frame}{Truth-Tracking with Non-Expert Sources}
    \begin{nitemize}
        \item AGM revision focusses on rationality, not on truth-tracking
        \item we adapt previous belief change problem with notions of
              truth-tracking
        \item see what can be learned in principle with non-experts: there are
              fundamental limits
          \item look at whether rationality is compatible with truth-tracking
              (yes)
    \end{nitemize}
\end{frame}

\end{document}
